%%
%% Copyright (c) 2018-2020 Weitian LI <wt@liwt.net>
%% CC BY 4.0 License
%%
%% Created: 2018-04-11
%%

% Chinese version
\documentclass[zh]{resume}


\name{维天}{李}

\keywords{BSD, Linux, Programming, Python, C, Shell, DevOps, SysAdmin}

% \tagline{\texorpdfstring{\icon{\faBinoculars} }{}<position-to-look-for>}
% \tagline{<current-position>}

% \photo[
%   shape=<circular|square>,    % default is circular
%   position=<left|right>,      % default is left
% ]{<width>}{<filename>}
% Example:
\photo[
shape=square,
position=right,
]{5.5em}{avatar.png}

\profile{
  \mobile{132-6262-0332}
  \email{liweitianux@live.com}
  \github{liweitianux} \\
  \university{上海交通大学}
  \degree{物理学 \textbullet 博士}
  \birthday{1991-09-26}
  \address{上海}
  % Custom information:
  % \icontext{<icon>}{<text>}
  % \iconlink{<icon>}{<link>}{<text>}
}

\begin{document}
\makeheader

%======================================================================
% Summary & Objectives
%======================================================================
\begin{abstract}
物理学专业（射电天文方向）直博研究生，有扎实的物理、数学与统计学基础，
擅长数据建模与分析，热衷计算机和网络技术，
有 10 年的 Linux 和 BSD 使用经验，熟练掌握 Shell、Python 和 C 语言编程。
积极实践自由开源精神，
在 \link{https://github.com/liweitianux}{GitHub} 上分享多个项目，
是 \link{https://www.dragonflybsd.org}{DragonFly BSD} 操作系统的开发者，
并积极参与其他多个开源项目。
\end{abstract}

%======================================================================
\sectionTitle{教育背景}{\faGraduationCap}
%======================================================================
\begin{educations}
  \education%
    {2013.09}%
    [2019.09]%
    {上海交通大学}%
    {物理与天文学院}%
    {物理学}%
    {博士}

  \separator{0.5ex}
  \education%
    {2009.09}%
    [2013.06]%
    {上海交通大学}%
    {物理与天文系}%
    {应用物理学}%
    {学士}
\end{educations}

%======================================================================
\sectionTitle{学术成果}{\faAtom}
%======================================================================
\begin{itemize}
  \item 开发低频射电天空图像模拟软件：
    \link{https://github.com/liweitianux/fg21sim}{\texttt{FG21sim}}
  \item 开发程序实现~X~射线天文观测数据的半自动化分析：
    \link{https://github.com/liweitianux/chandra-acis-analysis}{\texttt{chandra-acis-analysis}}
  \item 利用卷积去噪自动编码器（CDAE）在频率维度分离微弱的宇宙再电离（EoR）信号
  \item 利用卷积神经网络（CNN）对 FIRST 巡天的射电星系图像根据形态特征进行分类
  \item 显著改进星系团射电晕的建模，并考虑低频干涉阵列的复杂仪器效应
  \item 改进~X~射线光谱拟合的背景成分建模，获到更准确可靠的拟合结果
  \item 发表 2 篇第一作者以及 8 篇合作者 SCI 论文
\end{itemize}

%======================================================================
\sectionTitle{科研经历}{\faFlask}
%======================================================================
\begin{itemize}
  \item \link{https://github.com/liweitianux/atoolbox}{\texttt{atoolbox}}:
    (Python, Shell)
    多年来累积的各种工具，帮助管理系统、执行常用任务、分析天文数据等
  \item \link{https://github.com/liweitianux/dfly-update}{\texttt{dfly-update}}:
    (Shell)
    DragonFly BSD 系统更新程序
  \item \link{https://github.com/liweitianux/openrcs}{\texttt{openrcs}}:
    (C)
    改进 \texttt{OpenBSD RCS}，使其与 \texttt{GNU RCS} 足够兼容
  \item \link{https://github.com/liweitianux/fg21sim}{\texttt{fg21sim}}:
    (Python)
    模拟低频射电天空图像
  \item \link{https://github.com/liweitianux/cdae-eor}{\texttt{cdae-eor}}:
    (Python, Keras)
    使用卷积去噪自动编码器（CDAE）分离宇宙再电离（EoR）信号
  \item \link{https://github.com/liweitianux/chandra-acis-analysis}{\texttt{chandra-acis-analysis}}:
    (Python, Shell, Tcl)
    X~射线天文观测数据的半自动化分析程序
  \item \link{https://github.com/liweitianux/resume}{\texttt{resume}}:
    (\LaTeX)
    \emph{此简历}的模板和源文件
\end{itemize}


\begin{experiences}
  \experience%
    [2018.04]%
    {2018.08}%
    {数据工程师 @ 上海领脉网络科技（初创公司）}%
    [\begin{itemize}
      \item 从 Amazon 网页搜索并挖取商品与广告信息
        （Python, Requests, BeautifulSoup）
      \item 配置 Airflow 服务器和数据库等基础设施，
        定期从 Amazon 获取产品销售与广告投放等数据
      \item 开发网站（Flask, jQuery），帮助客户优化 Amazon 广告投放
    \end{itemize}]

  \separator{0.5ex}
  \experience%
    [2013.07]%
    {2013.09}%
    {网站开发 @ 97 随访（初创公司）}%
    [\begin{itemize}
      \item 后端开发（Django），完成用户注册、数据存储和搜索等功能
      \item 前端开发（jQuery, AJAX），对患者各项指标随时间的变化进行可视化
    \end{itemize}]
\end{experiences}

%======================================================================
\sectionTitle{科研技能}{\faWrench}
%======================================================================
\begin{competences}
  \comptence{编程与软件}{%
    Python（NumPy, SciPy, Astropy, Pandas）,%
    \enspace C/C++, Fortran, MATLAB, R, CUDA
  }
  \comptence{数据分析与机器学习}{%
    贝叶斯推断, MCMC, 深度学习（PyTorch, TensorFlow）,%
    \enspace 信号 / 图像处理, 统计建模
  }
  \comptence{射电天文}{%
    干涉阵数据校准与成像（CASA, MIRIAD）,%
    \enspace 天空图像模拟, 源提取（SExtractor）
  }
  \comptence{高性能计算}{%
    MPI, OpenMP, SLURM, GPU 编程, 集群管理
  }
  \comptence{科研工具}{%
    \LaTeX, Git, Docker, Jupyter, GNUPlot, FITS 数据处理
  }
  \comptence{\icon{\faLanguage} 语言}{
    \textbf{英语} --- 流利（读写听说）；%
    \enspace \textbf{中文} --- 母语
  }
\end{competences}

\begin{itemize}
  \item 具有丰富的射电干涉阵数据处理经验，
    处理过来自 SKA、LOFAR、VLA 和 Chandra 的观测数据
  \item 开发了端到端仿真管线，用于预测未来射电望远镜的性能和巡天策略
  \item 熟练使用贝叶斯统计方法进行天体物理参数估计和模型选择
  \item 搭建并维护了一套 4 节点的高性能计算集群，
    用于射电天空图像仿真和数据分析
  \item 设计并训练了卷积神经网络（CNN），
    用于射电星系形态分类
\end{itemize}

%======================================================================
\sectionTitle{荣誉奖项}{\faTrophy}
%======================================================================
\begin{itemize}
  \item 研究生国家奖学金（2017, 2018）
  \item 上海交通大学优秀学业奖学金（2018）
  \item 博士研究生一等奖助学金（2016, 2017）
  \item 中国天文学会年会优秀报告奖（2018）
  \item 国际天文学联合会会议旅行资助（2019）
\end{itemize}

\end{document}
